\section{Bonus: interpolation and Fourier inversion}
\label{sec:fourier-interpolation}

M\"obius inversion recovers exact contributions from cumulative sums.
Interpolation recovers a polynomial from its values. Both invert a
linear transformation: the observations are known linear combinations
of the data we want. Fourier inversion belongs to the same family.
At roots of unity, its connection with interpolation is exact.

\subsection{Evaluate a polynomial on a circle}

For this section, work over the complex numbers. Write $i^2=-1$ and
choose a \emph{primitive $n$th root of unity} $\omega$: its powers
$1,\omega,\ldots,\omega^{n-1}$ are distinct and $\omega^n=1$.
One choice is $\omega=\cos(2\pi/n)+i\sin(2\pi/n)$.
For $n=4$, we can use $\omega=i$, giving $1,i,-1,-i$.

Let
\[
 f(X)=\sum_{k=0}^{n-1}c_kX^k,
 \qquad y_j=f(\omega^j)=\sum_{k=0}^{n-1}c_k\omega^{jk}
 \quad(0\le j<n).
\]
With this choice of sign, the map $(c_k)\mapsto(y_j)$ is the
\emph{discrete Fourier transform}. Its matrix is the Vandermonde
matrix $(\omega^{jk})_{j,k}$: each row evaluates the monomials
$1,X,\ldots,X^{n-1}$ at one root of unity.
Some conventions put a minus sign in the forward transform; then
the inverse uses the opposite sign as well.

\begin{lemma}[Orthogonality of roots of unity]
For an integer $m$,
\[
 \sum_{j=0}^{n-1}\omega^{jm}
 =\begin{cases}n,&n\mid m,\\0,&n\nmid m.\end{cases}
\]
\end{lemma}
\begin{proof}
If $n\mid m$, every term is $1$. Otherwise $q=\omega^m\ne1$
and $q^n=1$, so the geometric sum is $(q^n-1)/(q-1)=0$.
\end{proof}

The functions $j\mapsto\omega^{jk}$ are the \emph{characters} of
the cyclic index set: addition of indices becomes multiplication
of their values. Orthogonality says that different characters
cancel when we sum over a complete cycle.

\subsection{Lagrange selectors are the inverse Fourier transform}

Recall the interpolation selectors from Chapter~4: the polynomial
$\ell_j$ of degree less than $n$ is $1$ at the $j$th sample point
and $0$ at the others. At roots of unity it has the particularly
simple form
\[
 \ell_j(X)=\frac1n\sum_{k=0}^{n-1}\omega^{-jk}X^k.
\]
Indeed, $\ell_j(\omega^r)=n^{-1}\sum_k\omega^{(r-j)k}=\delta_{jr}$
by orthogonality. These are exactly the Lagrange basis polynomials:
the degree bound and their values characterize them uniquely.
Interpolation therefore gives
\[
 f(X)=\sum_{j=0}^{n-1}y_j\ell_j(X)
 =\sum_{k=0}^{n-1}
   \left(\frac1n\sum_{j=0}^{n-1}y_j\omega^{-jk}\right)X^k.
\]
Comparing coefficients proves the inversion formula
\[
 \boxed{\displaystyle c_k=\frac1n\sum_{j=0}^{n-1}y_j\omega^{-jk}.}
\]
Thus interpolation at the roots of unity \emph{is} discrete Fourier
inversion. The factor $1/n$ comes from the surviving sum of $n$ ones.
The degree restriction matters: values at these points determine a
general polynomial only modulo $X^n-1$.

\begin{example}[Four samples recover four coefficients]
For $f(X)=1+2X+3X^2+4X^3$ and $\omega=i$, the samples are
\[
 (y_0,y_1,y_2,y_3)=(10,-2-2i,-2,-2+2i).
\]
The inverse formula becomes
\[
 \begin{aligned}
 c_0&=(y_0+y_1+y_2+y_3)/4=1,\\
 c_1&=(y_0-iy_1-y_2+iy_3)/4=2,\\
 c_2&=(y_0-y_1+y_2-y_3)/4=3,\\
 c_3&=(y_0+iy_1-y_2-iy_3)/4=4.
 \end{aligned}
\]
Changing the sign of the root in the inverse, as well as dividing
by $4$, is essential to recover the coefficients in the stated order.
\end{example}

\subsection{CRT explains multiplication in Fourier coordinates}

The factorization $X^n-1=\prod_{j=0}^{n-1}(X-\omega^j)$ has distinct
linear factors. Polynomial CRT identifies a remainder modulo $X^n-1$
with its $n$ values at these roots. The selectors $\ell_j$ reconstruct
that remainder, exactly as integer CRT selectors reconstruct residues.

If $f(X)=\sum c_kX^k$ and $g(X)=\sum d_kX^k$, their product modulo
$X^n-1$ has coefficients
\[
 (c\star d)_r=\sum_{k=0}^{n-1}c_kd_{\,r-k\bmod n}.
\]
This is \emph{cyclic convolution}: exponents wrap around because
$X^n=1$. Evaluation takes a product to the product of its values, so
\[
 \widehat{c\star d}_j=\widehat c_j\,\widehat d_j.
\]
Fourier coordinates turn cyclic convolution into pointwise multiplication.
Dirichlet convolution instead combines factor pairs $ab=n$, while
incidence convolution sums over intermediate elements of an interval.
Each convolution reflects the structure of its indexing set.

\subsection{What the three inversions share}

In each setting, compose the forward and inverse kernels and all
unwanted contributions cancel:
\[
 \begin{array}{ll}
 \text{M\"obius:}&\displaystyle
   \sum_{x\le z\le y}\mu_P(x,z)=\delta_{xy},\\[6pt]
 \text{Interpolation:}&\ell_j(x_r)=\delta_{jr},\\[4pt]
 \text{Fourier:}&\displaystyle
   \frac1n\sum_{j=0}^{n-1}\omega^{j(r-k)}=\delta_{rk}
   \quad(0\le r,k<n).
 \end{array}
\]
M\"obius coefficients undo accumulation along an order; Lagrange
polynomials undo evaluation at distinct points; Fourier coefficients
are the explicit interpolation inverse at roots of unity.
This is the common principle behind the three constructions:
choose coefficients that isolate one contribution and cancel the rest.

\begin{labbox}[breakable=false]
No analysis or limiting process was used in the Fourier argument.
The same formulas work over a field containing a primitive $n$th root
of unity, provided $n$ is invertible there. For example, in $\F_5$,
$2$ has order $4$, so the four-point transform uses $1,2,4,3$
and $1/4=4$. Notebook~6 checks both this example and the complex
four-point example with exact arithmetic. If the required roots are
missing, an extension field may supply them; that is a later use
of the arithmetic developed in Appendix~\ref{app:finite-fields}.
\end{labbox}
